\documentclass[11pt]{article}

\usepackage[margin=1in]{geometry}
\usepackage{amsmath,amssymb,amsthm,mathtools}
\usepackage{bm}
\usepackage{microtype}
\usepackage{enumitem}
\usepackage{booktabs}
\usepackage{array}
\usepackage{xcolor}
\usepackage{hyperref}
\usepackage[nameinlink,noabbrev]{cleveref}

\hypersetup{
  colorlinks=true,
  linkcolor=blue!60!black,
  citecolor=blue!60!black,
  urlcolor=blue!60!black
}

\theoremstyle{definition}
\newtheorem{definition}{Definition}[section]
\newtheorem{assumption}{Assumption}[section]
\theoremstyle{plain}
\newtheorem{theorem}[definition]{Theorem}
\newtheorem{lemma}[definition]{Lemma}
\newtheorem{proposition}[definition]{Proposition}
\newtheorem{corollary}[definition]{Corollary}
\theoremstyle{remark}
\newtheorem{remark}[definition]{Remark}
\crefname{definition}{definition}{definitions}
\Crefname{definition}{Definition}{Definitions}
\crefname{assumption}{assumption}{assumptions}
\Crefname{assumption}{Assumption}{Assumptions}
\crefname{theorem}{theorem}{theorems}
\Crefname{theorem}{Theorem}{Theorems}
\crefname{lemma}{lemma}{lemmas}
\Crefname{lemma}{Lemma}{Lemmas}
\crefname{proposition}{proposition}{propositions}
\Crefname{proposition}{Proposition}{Propositions}
\crefname{corollary}{corollary}{corollaries}
\Crefname{corollary}{Corollary}{Corollaries}
\crefname{remark}{remark}{remarks}
\Crefname{remark}{Remark}{Remarks}

\newcommand{\cH}{\mathcal H}
\newcommand{\cA}{\mathcal A}
\newcommand{\cD}{\mathcal D}
\newcommand{\cB}{\mathcal B}
\newcommand{\cC}{\mathcal C}
\newcommand{\cL}{\mathcal L}
\newcommand{\cK}{\mathcal K}
\newcommand{\cR}{\mathcal R}
\newcommand{\cG}{\mathcal G}
\newcommand{\Ztwo}{\mathbb Z_2}
\newcommand{\Sys}{\mathrm{Sys}}
\newcommand{\Time}{\mathrm{Time}}
\newcommand{\RG}{\mathrm{RG}}
\newcommand{\Tr}{\mathrm{Tr}}
\newcommand{\diag}{\mathrm{diag}}
\newcommand{\ad}{\mathrm{ad}}
\newcommand{\id}{\mathrm{id}}
\newcommand{\norm}[1]{\left\lVert #1\right\rVert}
\newcommand{\abs}[1]{\left\lvert #1\right\rvert}
\newcommand{\ip}[2]{\left\langle #1,#2\right\rangle}

\title{Completion-Sector Split Model Note}
\author{Internal technical note}
\date{\today}

\begin{document}
\maketitle

\begin{abstract}
This wrapper compiles the split-model section containing the completion-sector assumptions,
leak-law organization, threshold analysis, and the theorem-level nondegeneracy clause for the
normalized seam coefficient.
It serves as a stable `xr` target for those labels.
\end{abstract}

\tableofcontents

\IfFileExists{completion_sector_split_model_body.tex}{%
  \section{Completion-sector leak model and split criterion}
\label{sec:resolution}

This section separates what is directly derived from UOM from the additional hypotheses required to turn those ingredients
into a concrete split-driving leak model. The new organizing principle is:
the primary hidden variable is a \emph{signed odd split amplitude} \(a_k\), while the positive stock
\(S_k:=a_k^2\) is the quantity that can enter scalar observables and leak laws.
We reserve \(R_k^{\rm abs}\) for the remaining curvature budget and
\[
\mathfrak r_k:=\min_{\Lambda\ge \Lambda_k} f_k(\Lambda)
\]
for the minimized one-wall residual.

\subsection{Odd split mode, output parity, and the rank-one normal form}
\label{subsec:leak-solution}

The only ``leak'' that directly competes with wall work in the \(\Lambda\)-solver is the open-system modification of the
wall energy balance allowed explicitly in UOM's SK/HJ wall equation:
\begin{equation}
\Delta E_{{\rm geom},k}\ =\ J_{{\rm wall},k}\ +\ Q_{{\rm leak},k},
\qquad Q_{{\rm leak},k}\ge 0,
\label{eq:resolution-wall-leak}
\end{equation}
This is the welded wall-balance form used in the earlier summability analysis.
This is a \emph{receiver-side energy complement} term. It must not be conflated with UOM's information-theoretic
bridge/leak instrument from the UOM main manuscript.

\paragraph{What is already derived in UOM.}
Four ingredients are fixed by the manuscript:
\begin{enumerate}[leftmargin=2em]
\item \textbf{Wall balance.} The reduced one-wall description may contain a nonnegative complement \(Q_{\rm leak}\) as in
\eqref{eq:resolution-wall-leak}.
\item \textbf{Quadratic welded SK data.} In the Gaussian sector,
\[
W_\Lambda[J]=\frac12\int J\,(G_{R,\Lambda}+i\,\mathcal N_\Lambda)\,J+O(J^3),
\qquad \mathcal N_\Lambda\ge 0.
\]
\item \textbf{Healthy receiver.} The active KR receiver has vanishing welded RP block, so a receiver-tangent RP penalty is
not the fundamental home of split-driving oddness.
\item \textbf{Continued RP defect.} The analytic-continued covariance has a negative OS mode with
\[
\lambda_{\min}(\mu_{ps},\Lambda)
=
-\alpha(\Lambda)\mu_{ps}^2+O(\mu_{ps}^4),
\qquad
\alpha(\Lambda)\sim \kappa_d\,\Lambda^{d+2}H^{-d}.
\]
\end{enumerate}

\begin{assumption}[Split-mode observability]
\label{ass:wall-match}
Near split onset there exists a one-dimensional odd line \(L_k\) with local coordinate \(a_k\in\mathbb R\) such that the
split involution acts by
\[
a_k\longmapsto -a_k.
\]
The true two-wall advantage
\[
\Delta F_k:=\Delta\mathcal F_{2{\rm w}-1{\rm w},k}
\]
and the minimized one-wall residual
\[
\mathfrak r_k:=\min_{\Lambda\ge \Lambda_k} f_k(\Lambda)
\]
are smooth even scalar outputs of this same local mode.
Equivalently, the 2T \(\Gamma\)-odd descendant direction and the UOM continued negative mode \(u_-(\Lambda_k)\) are
assumed to represent the same active split mode up to a smooth local identification.
\end{assumption}

\begin{lemma}[Evenness of the minimized residual]
\label{lem:residual-even}
Suppose the pre-minimized one-wall residual is a smooth function \(f_k(\Lambda,a)\) satisfying
\[
f_k(\Lambda,-a)=f_k(\Lambda,a),
\]
and suppose there is a unique local minimizer \(\Lambda_\ast(a)\) near \(a=0\) such that
\[
\partial_\Lambda f_k(\Lambda_\ast(a),a)=0,
\qquad
\partial_\Lambda^2 f_k(\Lambda_\ast(0),0)>0.
\]
Then
\[
\mathfrak r_k(a):=f_k(\Lambda_\ast(a),a)
\]
is a smooth even function of \(a\).
\end{lemma}

\begin{proof}
By the implicit function theorem, \(\Lambda_\ast(a)\) is smooth near \(a=0\).
Since \(f_k(\Lambda,a)\) is even in \(a\), the stationarity equation for \(a\) and \(-a\) is identical; uniqueness of the
local minimizer gives
\[
\Lambda_\ast(-a)=\Lambda_\ast(a).
\]
Therefore
\[
\mathfrak r_k(-a)=f_k(\Lambda_\ast(-a),-a)=f_k(\Lambda_\ast(a),a)=\mathfrak r_k(a),
\]
so \(\mathfrak r_k\) is smooth and even.
\end{proof}

\begin{proposition}[Rank-one normal form near split onset]
\label{prop:rank-one-wall-match}
Under \cref{ass:wall-match}, the two scalar outputs admit local even expansions
\begin{equation}
\Delta F_k(a)
=
\mu_k^{(2{\rm w})}+h_k a^2+O(a^4),
\label{eq:DeltaF-detuned}
\end{equation}
\begin{equation}
\mathfrak r_k(a)
=
\mu_k^{(1{\rm w})}+m_k a^2+O(a^4),
\label{eq:rres-detuned}
\end{equation}
with \(h_k,m_k>0\) on the active split line.
\end{proposition}

\begin{proof}
Smooth even functions of one real variable have Taylor series containing only even powers.
The constants \(\mu_k^{(2{\rm w})}\) and \(\mu_k^{(1{\rm w})}\) are the baseline detunings of the two-wall and one-wall
channels. Positivity of \(h_k\) and \(m_k\) expresses that the selected mode is active in both outputs.
\end{proof}

\begin{corollary}[Onset wall-match in rank one]
\label{cor:rank-one-wall-match}
On the onset manifold,
\[
\mu_k^{(1{\rm w})}=\mu_k^{(2{\rm w})}=0,
\]
so
\[
\Delta F_k=h_k a_k^2+O(a_k^4),
\qquad
\mathfrak r_k=m_k a_k^2+O(a_k^4).
\]
Eliminating \(a_k^2\) gives
\begin{equation}
\Delta F_k
=
\frac{h_k}{m_k}\,\mathfrak r_k+O(\mathfrak r_k^2).
\label{eq:rank-one-wall-match}
\end{equation}
\end{corollary}

\subsection{Low-rank generalization}
\label{subsec:alpha-derived}

If the active split sector is \(m\)-dimensional, with odd amplitude vector \(a_k\in\mathbb R^m\), parity gives quadratic
normal forms
\begin{equation}
\Delta F_k(a)
=
\mu_k^{(2{\rm w})}+a^\top H_k a+O(\|a\|^4),
\label{eq:DeltaF-lowrank}
\end{equation}
\begin{equation}
\mathfrak r_k(a)
=
\mu_k^{(1{\rm w})}+a^\top M_k a+O(\|a\|^4),
\label{eq:rres-lowrank}
\end{equation}
with \(H_k,M_k\succeq 0\) on the active split subspace.

\begin{assumption}[Low-rank comparability]
\label{ass:low-rank-comp}
On the active split subspace there exist \(c_k^-,c_k^+>0\) such that
\[
c_k^- M_k \preceq H_k \preceq c_k^+ M_k.
\]
\end{assumption}

\begin{proposition}[Low-rank wall-match comparability]
\label{prop:low-rank-wall-match}
Assume \cref{ass:wall-match,ass:low-rank-comp} and restrict to the onset manifold
\(\mu_k^{(1{\rm w})}=\mu_k^{(2{\rm w})}=0\).
Then for sufficiently small \(\mathfrak r_k\) there exists \(C_k>0\) such that
\begin{equation}
c_k^-\,\mathfrak r_k-C_k\mathfrak r_k^2
\le
\Delta F_k
\le
c_k^+\,\mathfrak r_k+C_k\mathfrak r_k^2.
\label{eq:low-rank-wall-match}
\end{equation}
\end{proposition}

\begin{proof}
On onset,
\[
\Delta F_k=a_k^\top H_k a_k+O(\|a_k\|^4),
\qquad
\mathfrak r_k=a_k^\top M_k a_k+O(\|a_k\|^4).
\]
The quadratic-form comparison gives
\[
c_k^- a_k^\top M_k a_k\le a_k^\top H_k a_k\le c_k^+ a_k^\top M_k a_k.
\]
Since \(\mathfrak r_k\) is itself quadratic in \(a_k\) to leading order, the quartic remainders are \(O(\mathfrak r_k^2)\),
which yields \eqref{eq:low-rank-wall-match}.
\end{proof}

\subsection{Gaussian realization of the odd mode and the leak stock law}
\label{subsec:B-gaussian}

The theorem path above does \emph{not} rely on a KMS-subtracted matrix.
In particular, we do not use
\[
\Sigma_{B,k}:=\bigl(C_{B,k}-C^{\rm KMS}_{\Lambda_k}\bigr)+m_{B,k}m_{B,k}^\top
\]
as a positive object, because it need not be positive semidefinite.
When a safe matrix observable is needed in the realization layer, the natural PSD quantity is
\begin{equation}
M_{B,k}:=C_{B,k}+m_{B,k}m_{B,k}^\top\succeq 0.
\label{eq:MB-def}
\end{equation}

\begin{assumption}[Gaussian completion realization]
\label{ass:gaussian-completion}
For each admissible cutoff \(\Lambda\) in the welded regime, the quadratic influence functional \(W_\Lambda[J]\) admits a
Gaussian dilation, unique up to symplectic equivalence, by a bosonic completion sector \(B_\Lambda\) with canonical
quadratures \(X_\Lambda\) and positive quadratic Hamiltonian
\[
H_{B,\Lambda}=\frac12 X_\Lambda^\top h_\Lambda X_\Lambda,
\qquad h_\Lambda\succeq 0.
\]
\end{assumption}

\begin{assumption}[Dominant odd-mode projection]
\label{ass:export-selection}
Inside the Gaussian completion of \cref{ass:gaussian-completion}, there exists a smooth unit vector
\[
u_k:=u_-(\Lambda_k)
\]
spanning the active odd line and separated from the orthogonal sector by a local spectral gap.
We write
\[
P_k^\perp:=I-u_k u_k^\top
\]
for the orthogonal projector.
\end{assumption}

\begin{assumption}[Classical-control gate]
\label{ass:classical-gate}
There exists a scalar gate \(g_{\rm cl}(k)\in[0,1]\), determined by the accessible classicality defect
\(\mathrm{cl}_k\) or \(\mathrm{Off}_k\), such that \(g_{\rm cl}(k)\ll 1\) at a fresh-aeon start and
\(g_{\rm cl}(k)\to 1\) late in the aeon.
\end{assumption}

\begin{proposition}[Projected Gaussian affine dynamics]
\label{prop:projected-affine}
Under \cref{ass:gaussian-completion,ass:export-selection}, one may choose a finite effective completion-mode basis in which
the completion quadratures obey a linear-Gaussian recursion
\begin{equation}
X_{k+1}=A_k X_k+B_k j_k^{\rm ctl}+\nu_k,
\label{eq:X-recursion}
\end{equation}
with
\[
\mathbb E[\nu_k\mid\mathcal F_k]=0,
\qquad
\mathbb E[\nu_k\nu_k^\top\mid\mathcal F_k]=Q_k\succeq 0.
\]
Consequently the PSD second moment
\[
M_{B,k}=C_{B,k}+m_{B,k}m_{B,k}^\top
\]
obeys
\begin{equation}
M_{B,k+1}=A_k M_{B,k}A_k^\top+Q_k+\text{drift terms}.
\label{eq:MB-recursion}
\end{equation}
Projecting onto the odd line via
\[
a_k:=u_k^\top X_k
\]
gives the scalar affine law
\begin{equation}
a_{k+1}
=
\rho_k a_k+\beta_k+\xi_k+\varepsilon_k^{\rm mix},
\label{eq:B-recursion-general}
\end{equation}
where
\[
\rho_k:=u_{k+1}^\top A_k u_k,\qquad
\beta_k:=u_{k+1}^\top B_k j_k^{\rm ctl},
\]
\[
\xi_k:=u_{k+1}^\top \nu_k,\qquad
\varepsilon_k^{\rm mix}:=u_{k+1}^\top A_k P_k^\perp X_k.
\]
Moreover,
\[
\mathbb E[\xi_k\mid\mathcal F_k]=0,
\qquad
\mathbb E[\xi_k^2\mid\mathcal F_k]=u_{k+1}^\top Q_k u_{k+1}\ge 0.
\]
If \cref{ass:classical-gate} also holds and the welded shape family is fixed, then
\begin{equation}
\mathbb E[\xi_k^2\mid\mathcal F_k]
=
\eta_k\,g_{\rm cl}(k)\,\Delta I_k+O(\Delta I_k^2,\delta\Lambda_k^2),
\label{eq:Zk-linear}
\end{equation}
with \(\eta_k\ge 0\).
\end{proposition}

\begin{proof}
The quadratic SK influence functional determines a Gaussian CPTP map on the finite active completion-mode sector, hence a
linear-Gaussian state-space realization of the form \eqref{eq:X-recursion}; \eqref{eq:MB-recursion} is the standard update
for the PSD second moment.
Decompose \(X_k=u_k a_k+P_k^\perp X_k\), apply \(\ip{u_{k+1}}{\cdot}\) to \eqref{eq:X-recursion}, and collect the projected
transport, deterministic drive, noise, and orthogonal-mode contamination to obtain \eqref{eq:B-recursion-general}.
The mean-zero and positive-variance statements for \(\xi_k\) follow directly from the projected Gaussian noise covariance
\(
u_{k+1}^\top Q_k u_{k+1}\ge 0
\).
Finally, at fixed welded shape family one has \(A_{{\rm eff},k}^2\propto \Delta I_k\), so the fresh covariance injection
into the odd mode is linear in the current payload at leading order; the gate \(g_{\rm cl}(k)\) suppresses this injection
near a fresh-aeon start and yields \eqref{eq:Zk-linear}.
\end{proof}

\begin{assumption}[Centered odd preparation]
\label{ass:centered-odd}
Near split onset, the active odd mode is prepared either branch-symmetrically or in a squeeze-dominated regime, so the
projected deterministic displacement vanishes and orthogonal-mode contamination is higher order:
\[
\beta_k=0,
\qquad
\varepsilon_k^{\rm mix}=O_2.
\]
\end{assumption}

\begin{corollary}[Centered odd-mode law]
\label{cor:centered-odd-mode}
Under
\cref{ass:gaussian-completion,ass:export-selection,ass:classical-gate,ass:centered-odd},
the projected affine law \eqref{eq:B-recursion-general} reduces to
\begin{equation}
a_{k+1}=\rho_k a_k+\xi_k+O_2,
\label{eq:B-recursion}
\end{equation}
with \(\mathbb E[\xi_k\mid\mathcal F_k]=0\) and variance scaling \eqref{eq:Zk-linear}.
\end{corollary}

\paragraph{Stochastic/Gaussian realization layer.}
At theorem level the positive stock is simply
\[
S_k:=a_k^2.
\]
The general projected Gaussian dynamics need not be centered.
Under \cref{ass:centered-odd}, however, the physically relevant late-time regime is the one in which the odd mode carries no
leading deterministic displacement and the fresh effect enters through its variance.
In that centered specialization it is natural to use the conditional second moment
\[
S_k:=\mathbb E[a_k^2\mid\mathcal F_k].
\]

\begin{proposition}[Odd-mode stock law]
\label{prop:odd-mode-stock}
Under
\cref{ass:gaussian-completion,ass:export-selection,ass:classical-gate,ass:centered-odd},
the conditional stock
\[
S_k:=\mathbb E[a_k^2\mid\mathcal F_k]
\]
satisfies
\begin{equation}
S_{k+1}
=
t_k\,S_k+\eta_k\,g_{\rm cl}(k)\,\Delta I_k
+O(S_k\Delta I_k,\Delta I_k^2,\delta\Lambda_k^2),
\qquad t_k=\rho_k^2.
\label{eq:Sk-update}
\end{equation}
\end{proposition}

\begin{proof}
Square \eqref{eq:B-recursion}:
\[
a_{k+1}^2=\rho_k^2 a_k^2+2\rho_k a_k\xi_k+\xi_k^2+O_2.
\]
Taking conditional expectation and using \(\mathbb E[\xi_k\mid\mathcal F_k]=0\) removes the mixed term.
Substituting \eqref{eq:Zk-linear} gives \eqref{eq:Sk-update}.
\end{proof}

\begin{assumption}[UV comparability of energetic stiffness]
\label{ass:uv-match}
There exist constants \(c_-,c_+>0\) and a monotone increasing energetic stiffness \(\widetilde\alpha(\Lambda)\) of the
active odd mode such that, for all sufficiently large \(\Lambda\),
\begin{equation}
c_-\,\alpha(\Lambda)\le \widetilde\alpha(\Lambda)\le c_+\,\alpha(\Lambda).
\label{eq:alpha-comparable}
\end{equation}
\end{assumption}

This is the precise place where the note passes from what is derived in UOM to a physical hypothesis.
The coefficient \(\alpha(\Lambda)\) is derived from the continued OS problem; the statement that the energetic stiffness of
the split mode is comparable to it is not yet a theorem.

With the stock \(S_k\) as the actual state variable, the positive split-driving export flux is modeled by
\begin{equation}
Q_{{\rm leak},k}^{\rm model}
=
\widetilde\alpha'(\Lambda_k)\,\delta\Lambda_k\,S_k
\;+\;
\widetilde\alpha(\Lambda_k)\,\eta_k\,g_{\rm cl}(k)\,\Delta I_k
\;+\; O_2,
\label{eq:Qleak-effective}
\end{equation}
where \(\delta\Lambda_k:=\Lambda_{k+1}-\Lambda_k\).
The first term is the transport cost of carrying already-prepared odd stock to larger cutoff; the second is the fresh
injection due to the current payload.
Under \cref{ass:uv-match}, this gives the scaling comparison
\begin{equation}
Q_{{\rm leak},k}^{\rm model}
\asymp
\alpha'(\Lambda_k)\,\delta\Lambda_k\,S_k
\;+\;
\alpha(\Lambda_k)\,\eta_k\,g_{\rm cl}(k)\,\Delta I_k.
\label{eq:Qleak-alpha}
\end{equation}

\paragraph{Current RP proxy reinterpreted.}
The quantity
\[
Q_{\rm rg}=\chi J_{\rm RP}/\beta_{\rm dS}
\]
therefore remains useful only as a phenomenological probe. It may correlate with \(a_k^2\) or with the stock \(S_k\), but
it is not the fundamental leak law because the healthy receiver itself has no quadratic RP block on the welded family.

\subsection{Approximate split condition from the odd-stock leak law}
\label{subsec:split-threshold}

Let \(R_k^{\rm abs}:=|R_k|\). In the micro-step regime the geometry-side surrogate gives
\[
\Delta E_{{\rm geom},k}
\approx
\frac{\zeta\,(A+c_\Lambda \Lambda_k^2)}{2\sqrt{12\,R_k^{\rm abs}}}\;\Theta_k
\;-\;
\frac{2\zeta\,c_\Lambda \Lambda_k\sqrt{R_k^{\rm abs}}}{\sqrt{12}}\;\delta\Lambda,
\]
while
\[
J_{{\rm wall},k}\approx \mu_{\rm RG}(\Lambda_k)\,\delta\Lambda,
\qquad
\mu_{\rm RG}(\Lambda)\simeq \mu_0+\mu_2\Lambda^2>0.
\]
Hence the local one-wall threshold remains
\begin{equation}
Q_{{\rm leak},k}
\gtrsim
\frac{\zeta\,(A+c_\Lambda \Lambda_k^2)}{2\sqrt{12\,R_k^{\rm abs}}}\;\Theta_k.
\label{eq:split-threshold-base}
\end{equation}
Since
\[
\Theta_k
=
K(\Lambda_k)\,\Delta I_k
\sim
K_0\,\Lambda_k^{-p}\,e^{-2d_{\rm sep}(\Lambda_k)\Lambda_k}\,\Delta I_k,
\qquad p\in\{10,14\},
\]
substituting \eqref{eq:Qleak-effective} gives
\begin{equation}
\widetilde\alpha'(\Lambda_k)\,\delta\Lambda_k\,S_k
\;+\;
\widetilde\alpha(\Lambda_k)\,\eta_k\,g_{\rm cl}(k)\,\Delta I_k
\gtrsim
\frac{\zeta\,(A+c_\Lambda \Lambda_k^2)}{2\sqrt{12\,R_k^{\rm abs}}}\,
K(\Lambda_k)\,\Delta I_k.
\label{eq:split-threshold-effective}
\end{equation}
This is the corrected split threshold. It shows the two late-time branches transparently:
\begin{enumerate}[leftmargin=2em]
\item if \(S_k\) remains negligible, only the fresh term remains, and the exhaustion branch can survive;
\item if \(S_k\) has accumulated, the transport term can dominate even when \(\Delta I_k\) is already tiny.
\end{enumerate}
Under \cref{ass:uv-match} and \(\alpha(\Lambda)\sim \kappa_d\Lambda^{d+2}H^{-d}\), the transport term scales like
\[
\widetilde\alpha'(\Lambda_k)\,\delta\Lambda_k\,S_k
\asymp
\Lambda_k^{d+1}\,\delta\Lambda_k\,S_k,
\]
which is exactly the mechanism by which late split can beat the shrinking fresh payload.

\subsection{Nondegeneracy clause for the normalized seam coefficient}
\label{subsec:separate-c-nondegeneracy}

The split threshold \eqref{eq:split-threshold-effective} is a statement about when the late-time
odd stock can overcome the shrinking fresh payload.
It is \emph{not} by itself a theorem that the normalized seam coefficient \(c_{\Lambda_k}\) on the
canonical compressed welded LoG line is nonzero.
Any theorem proving \(c_{\Lambda_k}\neq 0\) must therefore come from a separate overlap or transport
hypothesis, not from the threshold inequality alone.
Downstream, this clause is absorbed into the delayed selected-line package rather than carried as a
separate package summand.

\begin{definition}[Unit-source seam overlap scalar and no-destructive-cancellation margin]
\label{def:unit-source-overlap-margin}
Fix an accepted tick \(k\).
Assume the canonical compressed welded LoG line is defined and let
\[
\tau^{\mathrm{unit}}_{\Lambda_k}
\]
be a detector-factorized unit welded odd source.
Define the raw unit-source seam overlap scalar by
\begin{equation}
\widetilde c_{\Lambda_k}
:=
\big\langle\!\big\langle
\Pi_{\mathrm{phys},\Lambda_k}\mathsf T_{\Lambda_k}\!\big[\tau^{\mathrm{unit}}_{\Lambda_k}\big],\,
\delta C^{\mathrm{LoG}}_{\Lambda_k}
\big\rangle\!\big\rangle_{C_{\ast,\Lambda_k}}.
\label{eq:unit-source-overlap-scalar}
\end{equation}
A \emph{no-destructive-cancellation margin} at \(\Lambda_k\) is any real number
\[
\mu_{\Lambda_k}>0
\]
such that
\begin{equation}
\abs{\widetilde c_{\Lambda_k}}\ge \mu_{\Lambda_k}.
\label{eq:no-destr-cancel-margin}
\end{equation}
On a connected accepted-branch segment \(I\), a \emph{branchwise no-destructive-cancellation
margin} is a constant \(\mu_I>0\) such that
\begin{equation}
\abs{\widetilde c_{\Lambda}}\ge \mu_I
\qquad
\forall\,\Lambda\in I.
\label{eq:no-destr-cancel-margin-branch}
\end{equation}
\end{definition}

\begin{theorem}[Nondegeneracy criterion for \texorpdfstring{\(c_{\Lambda_k}\)}{c_Lambda_k}]
\label{thm:separate-c-nondegeneracy}
Fix an accepted tick \(k\).
Assume:
\begin{enumerate}[leftmargin=2em]
\item a detector-factorized unit welded odd source at \(\Lambda_k\);
\item the unit-source seam-faithfulness package, namely:
  \begin{enumerate}[label=(\alph*),leftmargin=2em]
  \item the raw seam pairing factorizes as
    \begin{equation}
    \widetilde s_{\Lambda_k}(f,y)
    =
    \sigma_{\mathrm{fill}}(f)\,
    \widetilde c_{\Lambda_k}\,
    \overline{\mathrm{obs}}(y)
    \qquad
    \forall f,\forall y;
    \label{eq:unit-source-seam-factorization}
    \end{equation}
  \item whenever
    \(
    \widetilde c_{\Lambda_k}\neq 0
    \)
    and
    \(
    \overline{\mathrm{obs}}(y)\neq 0
    \),
    the compressed welded LoG line is transversely physical, dominant-line realization holds, and
    the normalized seam coefficient
    \begin{equation}
    c_{\Lambda_k}
    :=
    \kappa_{\Lambda_k}^{-1/2}\,\widetilde c_{\Lambda_k}
    \label{eq:unit-source-normalized-c}
    \end{equation}
    is defined;
  \end{enumerate}
\item a no-destructive-cancellation margin at \(\Lambda_k\) in the sense of
  \cref{def:unit-source-overlap-margin};
\item a witness \(y\) with
  \[
  \overline{\mathrm{obs}}(y)\neq 0.
  \]
\end{enumerate}
Then:
\begin{enumerate}[leftmargin=2em]
\item the raw seam pairing is nonzero for every filler tag \(f\):
  \[
  \widetilde s_{\Lambda_k}(f,y)\neq 0;
  \]
\item the canonical compressed welded LoG line is transversely physical at \(\Lambda_k\);
\item dominant-line realization holds at \(\Lambda_k\);
\item the normalized seam coefficient is defined and nonzero:
  \[
  c_{\Lambda_k}\neq 0.
  \]
\end{enumerate}
If, in addition, for some filler tag \(f\) the source \(\tau^{f,y}_{\Lambda_k}\) is representable on
the characterization layer and admits a source-compatible descendant realization, then the
corresponding descendant representative satisfies
\[
\rho^{\mathrm{dom}}_{\Lambda_k}\!\big(J^{f,y}_{\Lambda_k}\big)\neq 0.
\]
\end{theorem}

\begin{proof}
The margin \eqref{eq:no-destr-cancel-margin} implies
\(
\widetilde c_{\Lambda_k}\neq 0
\).
Together with
\(
\overline{\mathrm{obs}}(y)\neq 0
\),
the factorization \eqref{eq:unit-source-seam-factorization} gives
\(
\widetilde s_{\Lambda_k}(f,y)\neq 0
\)
for every filler tag \(f\), proving item (1).
Item (2), item (3), and the definition and nonvanishing of
\(
c_{\Lambda_k}
\)
are exactly the remaining conclusions of the assumed unit-source seam-faithfulness package under the
same nondegeneracy hypothesis.
The final sentence is the corresponding witness-to-descendant comparison step under representability
and source-compatible descendant realization.
\end{proof}

\begin{corollary}[Branchwise sign-stable nondegeneracy]
\label{cor:branchwise-c-nondegeneracy}
Let \(I\) be a connected accepted-branch segment.
Assume:
\begin{enumerate}[leftmargin=2em]
\item the branchwise normalized seam coefficient
  \[
  \Lambda\mapsto c_{\Lambda}
  \]
  is continuous on \(I\);
\item a branchwise no-destructive-cancellation margin on \(I\) in the sense of
  \cref{def:unit-source-overlap-margin}.
\end{enumerate}
Then
\[
c_{\Lambda}\neq 0
\qquad
\forall\,\Lambda\in I.
\]
In particular, the sign of \(c_{\Lambda}\) is constant on \(I\).
If there also exist a fixed LS witness \(y\) with
\(
\overline{\mathrm{obs}}(y)\neq 0
\)
and a branchwise source-compatible descendant realization satisfying
\[
\rho^{\mathrm{dom}}_{\Lambda}\!\big(J^{f,y}_{\Lambda}\big)
=
\sigma_{\mathrm{fill}}(f)\,
c_{\Lambda}\,
\overline{\mathrm{obs}}(y),
\]
then the sign of
\(
\rho^{\mathrm{dom}}_{\Lambda}(J^{f,y}_{\Lambda})
\)
is constant on \(I\) as well.
\end{corollary}

\begin{proof}
By \eqref{eq:no-destr-cancel-margin-branch}, the scalar
\(
\widetilde c_{\Lambda}
\)
never vanishes on \(I\).
Hence neither does
\(
c_{\Lambda}
\),
because the normalization only rescales the raw overlap scalar by a positive branchwise factor.
A continuous nonzero real-valued function on a connected segment cannot change sign, proving the
constancy of \(\operatorname{sgn}(c_{\Lambda})\).
The final statement follows immediately from the displayed branchwise comparison identity and the
fixed nonzero scalar \(\overline{\mathrm{obs}}(y)\).
\end{proof}

\begin{remark}[Logical separation from the split threshold]
\label{rem:separate-c-nondegeneracy}
The hypotheses of \cref{thm:separate-c-nondegeneracy,cor:branchwise-c-nondegeneracy} are
intentionally separate from \eqref{eq:split-threshold-effective}.
The threshold governs the emergence of split-side stock; the no-destructive-cancellation margin
governs overlap of the unit welded odd source with the canonical compressed dominant line.
Accordingly, the present note treats \(c_{\Lambda}\neq 0\) as the transport/overlap clause later
imported into the delayed selected-line package rather than as a consequence of the split
threshold.
\end{remark}
%
}{%
  \IfFileExists{UOM/completion_sector_split_model_body.tex}{%
    \section{Completion-sector leak model and split criterion}
\label{sec:resolution}

This section separates what is directly derived from UOM from the additional hypotheses required to turn those ingredients
into a concrete split-driving leak model. The new organizing principle is:
the primary hidden variable is a \emph{signed odd split amplitude} \(a_k\), while the positive stock
\(S_k:=a_k^2\) is the quantity that can enter scalar observables and leak laws.
We reserve \(R_k^{\rm abs}\) for the remaining curvature budget and
\[
\mathfrak r_k:=\min_{\Lambda\ge \Lambda_k} f_k(\Lambda)
\]
for the minimized one-wall residual.

\subsection{Odd split mode, output parity, and the rank-one normal form}
\label{subsec:leak-solution}

The only ``leak'' that directly competes with wall work in the \(\Lambda\)-solver is the open-system modification of the
wall energy balance allowed explicitly in UOM's SK/HJ wall equation:
\begin{equation}
\Delta E_{{\rm geom},k}\ =\ J_{{\rm wall},k}\ +\ Q_{{\rm leak},k},
\qquad Q_{{\rm leak},k}\ge 0,
\label{eq:resolution-wall-leak}
\end{equation}
This is the welded wall-balance form used in the earlier summability analysis.
This is a \emph{receiver-side energy complement} term. It must not be conflated with UOM's information-theoretic
bridge/leak instrument from the UOM main manuscript.

\paragraph{What is already derived in UOM.}
Four ingredients are fixed by the manuscript:
\begin{enumerate}[leftmargin=2em]
\item \textbf{Wall balance.} The reduced one-wall description may contain a nonnegative complement \(Q_{\rm leak}\) as in
\eqref{eq:resolution-wall-leak}.
\item \textbf{Quadratic welded SK data.} In the Gaussian sector,
\[
W_\Lambda[J]=\frac12\int J\,(G_{R,\Lambda}+i\,\mathcal N_\Lambda)\,J+O(J^3),
\qquad \mathcal N_\Lambda\ge 0.
\]
\item \textbf{Healthy receiver.} The active KR receiver has vanishing welded RP block, so a receiver-tangent RP penalty is
not the fundamental home of split-driving oddness.
\item \textbf{Continued RP defect.} The analytic-continued covariance has a negative OS mode with
\[
\lambda_{\min}(\mu_{ps},\Lambda)
=
-\alpha(\Lambda)\mu_{ps}^2+O(\mu_{ps}^4),
\qquad
\alpha(\Lambda)\sim \kappa_d\,\Lambda^{d+2}H^{-d}.
\]
\end{enumerate}

\begin{assumption}[Split-mode observability]
\label{ass:wall-match}
Near split onset there exists a one-dimensional odd line \(L_k\) with local coordinate \(a_k\in\mathbb R\) such that the
split involution acts by
\[
a_k\longmapsto -a_k.
\]
The true two-wall advantage
\[
\Delta F_k:=\Delta\mathcal F_{2{\rm w}-1{\rm w},k}
\]
and the minimized one-wall residual
\[
\mathfrak r_k:=\min_{\Lambda\ge \Lambda_k} f_k(\Lambda)
\]
are smooth even scalar outputs of this same local mode.
Equivalently, the 2T \(\Gamma\)-odd descendant direction and the UOM continued negative mode \(u_-(\Lambda_k)\) are
assumed to represent the same active split mode up to a smooth local identification.
\end{assumption}

\begin{lemma}[Evenness of the minimized residual]
\label{lem:residual-even}
Suppose the pre-minimized one-wall residual is a smooth function \(f_k(\Lambda,a)\) satisfying
\[
f_k(\Lambda,-a)=f_k(\Lambda,a),
\]
and suppose there is a unique local minimizer \(\Lambda_\ast(a)\) near \(a=0\) such that
\[
\partial_\Lambda f_k(\Lambda_\ast(a),a)=0,
\qquad
\partial_\Lambda^2 f_k(\Lambda_\ast(0),0)>0.
\]
Then
\[
\mathfrak r_k(a):=f_k(\Lambda_\ast(a),a)
\]
is a smooth even function of \(a\).
\end{lemma}

\begin{proof}
By the implicit function theorem, \(\Lambda_\ast(a)\) is smooth near \(a=0\).
Since \(f_k(\Lambda,a)\) is even in \(a\), the stationarity equation for \(a\) and \(-a\) is identical; uniqueness of the
local minimizer gives
\[
\Lambda_\ast(-a)=\Lambda_\ast(a).
\]
Therefore
\[
\mathfrak r_k(-a)=f_k(\Lambda_\ast(-a),-a)=f_k(\Lambda_\ast(a),a)=\mathfrak r_k(a),
\]
so \(\mathfrak r_k\) is smooth and even.
\end{proof}

\begin{proposition}[Rank-one normal form near split onset]
\label{prop:rank-one-wall-match}
Under \cref{ass:wall-match}, the two scalar outputs admit local even expansions
\begin{equation}
\Delta F_k(a)
=
\mu_k^{(2{\rm w})}+h_k a^2+O(a^4),
\label{eq:DeltaF-detuned}
\end{equation}
\begin{equation}
\mathfrak r_k(a)
=
\mu_k^{(1{\rm w})}+m_k a^2+O(a^4),
\label{eq:rres-detuned}
\end{equation}
with \(h_k,m_k>0\) on the active split line.
\end{proposition}

\begin{proof}
Smooth even functions of one real variable have Taylor series containing only even powers.
The constants \(\mu_k^{(2{\rm w})}\) and \(\mu_k^{(1{\rm w})}\) are the baseline detunings of the two-wall and one-wall
channels. Positivity of \(h_k\) and \(m_k\) expresses that the selected mode is active in both outputs.
\end{proof}

\begin{corollary}[Onset wall-match in rank one]
\label{cor:rank-one-wall-match}
On the onset manifold,
\[
\mu_k^{(1{\rm w})}=\mu_k^{(2{\rm w})}=0,
\]
so
\[
\Delta F_k=h_k a_k^2+O(a_k^4),
\qquad
\mathfrak r_k=m_k a_k^2+O(a_k^4).
\]
Eliminating \(a_k^2\) gives
\begin{equation}
\Delta F_k
=
\frac{h_k}{m_k}\,\mathfrak r_k+O(\mathfrak r_k^2).
\label{eq:rank-one-wall-match}
\end{equation}
\end{corollary}

\subsection{Low-rank generalization}
\label{subsec:alpha-derived}

If the active split sector is \(m\)-dimensional, with odd amplitude vector \(a_k\in\mathbb R^m\), parity gives quadratic
normal forms
\begin{equation}
\Delta F_k(a)
=
\mu_k^{(2{\rm w})}+a^\top H_k a+O(\|a\|^4),
\label{eq:DeltaF-lowrank}
\end{equation}
\begin{equation}
\mathfrak r_k(a)
=
\mu_k^{(1{\rm w})}+a^\top M_k a+O(\|a\|^4),
\label{eq:rres-lowrank}
\end{equation}
with \(H_k,M_k\succeq 0\) on the active split subspace.

\begin{assumption}[Low-rank comparability]
\label{ass:low-rank-comp}
On the active split subspace there exist \(c_k^-,c_k^+>0\) such that
\[
c_k^- M_k \preceq H_k \preceq c_k^+ M_k.
\]
\end{assumption}

\begin{proposition}[Low-rank wall-match comparability]
\label{prop:low-rank-wall-match}
Assume \cref{ass:wall-match,ass:low-rank-comp} and restrict to the onset manifold
\(\mu_k^{(1{\rm w})}=\mu_k^{(2{\rm w})}=0\).
Then for sufficiently small \(\mathfrak r_k\) there exists \(C_k>0\) such that
\begin{equation}
c_k^-\,\mathfrak r_k-C_k\mathfrak r_k^2
\le
\Delta F_k
\le
c_k^+\,\mathfrak r_k+C_k\mathfrak r_k^2.
\label{eq:low-rank-wall-match}
\end{equation}
\end{proposition}

\begin{proof}
On onset,
\[
\Delta F_k=a_k^\top H_k a_k+O(\|a_k\|^4),
\qquad
\mathfrak r_k=a_k^\top M_k a_k+O(\|a_k\|^4).
\]
The quadratic-form comparison gives
\[
c_k^- a_k^\top M_k a_k\le a_k^\top H_k a_k\le c_k^+ a_k^\top M_k a_k.
\]
Since \(\mathfrak r_k\) is itself quadratic in \(a_k\) to leading order, the quartic remainders are \(O(\mathfrak r_k^2)\),
which yields \eqref{eq:low-rank-wall-match}.
\end{proof}

\subsection{Gaussian realization of the odd mode and the leak stock law}
\label{subsec:B-gaussian}

The theorem path above does \emph{not} rely on a KMS-subtracted matrix.
In particular, we do not use
\[
\Sigma_{B,k}:=\bigl(C_{B,k}-C^{\rm KMS}_{\Lambda_k}\bigr)+m_{B,k}m_{B,k}^\top
\]
as a positive object, because it need not be positive semidefinite.
When a safe matrix observable is needed in the realization layer, the natural PSD quantity is
\begin{equation}
M_{B,k}:=C_{B,k}+m_{B,k}m_{B,k}^\top\succeq 0.
\label{eq:MB-def}
\end{equation}

\begin{assumption}[Gaussian completion realization]
\label{ass:gaussian-completion}
For each admissible cutoff \(\Lambda\) in the welded regime, the quadratic influence functional \(W_\Lambda[J]\) admits a
Gaussian dilation, unique up to symplectic equivalence, by a bosonic completion sector \(B_\Lambda\) with canonical
quadratures \(X_\Lambda\) and positive quadratic Hamiltonian
\[
H_{B,\Lambda}=\frac12 X_\Lambda^\top h_\Lambda X_\Lambda,
\qquad h_\Lambda\succeq 0.
\]
\end{assumption}

\begin{assumption}[Dominant odd-mode projection]
\label{ass:export-selection}
Inside the Gaussian completion of \cref{ass:gaussian-completion}, there exists a smooth unit vector
\[
u_k:=u_-(\Lambda_k)
\]
spanning the active odd line and separated from the orthogonal sector by a local spectral gap.
We write
\[
P_k^\perp:=I-u_k u_k^\top
\]
for the orthogonal projector.
\end{assumption}

\begin{assumption}[Classical-control gate]
\label{ass:classical-gate}
There exists a scalar gate \(g_{\rm cl}(k)\in[0,1]\), determined by the accessible classicality defect
\(\mathrm{cl}_k\) or \(\mathrm{Off}_k\), such that \(g_{\rm cl}(k)\ll 1\) at a fresh-aeon start and
\(g_{\rm cl}(k)\to 1\) late in the aeon.
\end{assumption}

\begin{proposition}[Projected Gaussian affine dynamics]
\label{prop:projected-affine}
Under \cref{ass:gaussian-completion,ass:export-selection}, one may choose a finite effective completion-mode basis in which
the completion quadratures obey a linear-Gaussian recursion
\begin{equation}
X_{k+1}=A_k X_k+B_k j_k^{\rm ctl}+\nu_k,
\label{eq:X-recursion}
\end{equation}
with
\[
\mathbb E[\nu_k\mid\mathcal F_k]=0,
\qquad
\mathbb E[\nu_k\nu_k^\top\mid\mathcal F_k]=Q_k\succeq 0.
\]
Consequently the PSD second moment
\[
M_{B,k}=C_{B,k}+m_{B,k}m_{B,k}^\top
\]
obeys
\begin{equation}
M_{B,k+1}=A_k M_{B,k}A_k^\top+Q_k+\text{drift terms}.
\label{eq:MB-recursion}
\end{equation}
Projecting onto the odd line via
\[
a_k:=u_k^\top X_k
\]
gives the scalar affine law
\begin{equation}
a_{k+1}
=
\rho_k a_k+\beta_k+\xi_k+\varepsilon_k^{\rm mix},
\label{eq:B-recursion-general}
\end{equation}
where
\[
\rho_k:=u_{k+1}^\top A_k u_k,\qquad
\beta_k:=u_{k+1}^\top B_k j_k^{\rm ctl},
\]
\[
\xi_k:=u_{k+1}^\top \nu_k,\qquad
\varepsilon_k^{\rm mix}:=u_{k+1}^\top A_k P_k^\perp X_k.
\]
Moreover,
\[
\mathbb E[\xi_k\mid\mathcal F_k]=0,
\qquad
\mathbb E[\xi_k^2\mid\mathcal F_k]=u_{k+1}^\top Q_k u_{k+1}\ge 0.
\]
If \cref{ass:classical-gate} also holds and the welded shape family is fixed, then
\begin{equation}
\mathbb E[\xi_k^2\mid\mathcal F_k]
=
\eta_k\,g_{\rm cl}(k)\,\Delta I_k+O(\Delta I_k^2,\delta\Lambda_k^2),
\label{eq:Zk-linear}
\end{equation}
with \(\eta_k\ge 0\).
\end{proposition}

\begin{proof}
The quadratic SK influence functional determines a Gaussian CPTP map on the finite active completion-mode sector, hence a
linear-Gaussian state-space realization of the form \eqref{eq:X-recursion}; \eqref{eq:MB-recursion} is the standard update
for the PSD second moment.
Decompose \(X_k=u_k a_k+P_k^\perp X_k\), apply \(\ip{u_{k+1}}{\cdot}\) to \eqref{eq:X-recursion}, and collect the projected
transport, deterministic drive, noise, and orthogonal-mode contamination to obtain \eqref{eq:B-recursion-general}.
The mean-zero and positive-variance statements for \(\xi_k\) follow directly from the projected Gaussian noise covariance
\(
u_{k+1}^\top Q_k u_{k+1}\ge 0
\).
Finally, at fixed welded shape family one has \(A_{{\rm eff},k}^2\propto \Delta I_k\), so the fresh covariance injection
into the odd mode is linear in the current payload at leading order; the gate \(g_{\rm cl}(k)\) suppresses this injection
near a fresh-aeon start and yields \eqref{eq:Zk-linear}.
\end{proof}

\begin{assumption}[Centered odd preparation]
\label{ass:centered-odd}
Near split onset, the active odd mode is prepared either branch-symmetrically or in a squeeze-dominated regime, so the
projected deterministic displacement vanishes and orthogonal-mode contamination is higher order:
\[
\beta_k=0,
\qquad
\varepsilon_k^{\rm mix}=O_2.
\]
\end{assumption}

\begin{corollary}[Centered odd-mode law]
\label{cor:centered-odd-mode}
Under
\cref{ass:gaussian-completion,ass:export-selection,ass:classical-gate,ass:centered-odd},
the projected affine law \eqref{eq:B-recursion-general} reduces to
\begin{equation}
a_{k+1}=\rho_k a_k+\xi_k+O_2,
\label{eq:B-recursion}
\end{equation}
with \(\mathbb E[\xi_k\mid\mathcal F_k]=0\) and variance scaling \eqref{eq:Zk-linear}.
\end{corollary}

\paragraph{Stochastic/Gaussian realization layer.}
At theorem level the positive stock is simply
\[
S_k:=a_k^2.
\]
The general projected Gaussian dynamics need not be centered.
Under \cref{ass:centered-odd}, however, the physically relevant late-time regime is the one in which the odd mode carries no
leading deterministic displacement and the fresh effect enters through its variance.
In that centered specialization it is natural to use the conditional second moment
\[
S_k:=\mathbb E[a_k^2\mid\mathcal F_k].
\]

\begin{proposition}[Odd-mode stock law]
\label{prop:odd-mode-stock}
Under
\cref{ass:gaussian-completion,ass:export-selection,ass:classical-gate,ass:centered-odd},
the conditional stock
\[
S_k:=\mathbb E[a_k^2\mid\mathcal F_k]
\]
satisfies
\begin{equation}
S_{k+1}
=
t_k\,S_k+\eta_k\,g_{\rm cl}(k)\,\Delta I_k
+O(S_k\Delta I_k,\Delta I_k^2,\delta\Lambda_k^2),
\qquad t_k=\rho_k^2.
\label{eq:Sk-update}
\end{equation}
\end{proposition}

\begin{proof}
Square \eqref{eq:B-recursion}:
\[
a_{k+1}^2=\rho_k^2 a_k^2+2\rho_k a_k\xi_k+\xi_k^2+O_2.
\]
Taking conditional expectation and using \(\mathbb E[\xi_k\mid\mathcal F_k]=0\) removes the mixed term.
Substituting \eqref{eq:Zk-linear} gives \eqref{eq:Sk-update}.
\end{proof}

\begin{assumption}[UV comparability of energetic stiffness]
\label{ass:uv-match}
There exist constants \(c_-,c_+>0\) and a monotone increasing energetic stiffness \(\widetilde\alpha(\Lambda)\) of the
active odd mode such that, for all sufficiently large \(\Lambda\),
\begin{equation}
c_-\,\alpha(\Lambda)\le \widetilde\alpha(\Lambda)\le c_+\,\alpha(\Lambda).
\label{eq:alpha-comparable}
\end{equation}
\end{assumption}

This is the precise place where the note passes from what is derived in UOM to a physical hypothesis.
The coefficient \(\alpha(\Lambda)\) is derived from the continued OS problem; the statement that the energetic stiffness of
the split mode is comparable to it is not yet a theorem.

With the stock \(S_k\) as the actual state variable, the positive split-driving export flux is modeled by
\begin{equation}
Q_{{\rm leak},k}^{\rm model}
=
\widetilde\alpha'(\Lambda_k)\,\delta\Lambda_k\,S_k
\;+\;
\widetilde\alpha(\Lambda_k)\,\eta_k\,g_{\rm cl}(k)\,\Delta I_k
\;+\; O_2,
\label{eq:Qleak-effective}
\end{equation}
where \(\delta\Lambda_k:=\Lambda_{k+1}-\Lambda_k\).
The first term is the transport cost of carrying already-prepared odd stock to larger cutoff; the second is the fresh
injection due to the current payload.
Under \cref{ass:uv-match}, this gives the scaling comparison
\begin{equation}
Q_{{\rm leak},k}^{\rm model}
\asymp
\alpha'(\Lambda_k)\,\delta\Lambda_k\,S_k
\;+\;
\alpha(\Lambda_k)\,\eta_k\,g_{\rm cl}(k)\,\Delta I_k.
\label{eq:Qleak-alpha}
\end{equation}

\paragraph{Current RP proxy reinterpreted.}
The quantity
\[
Q_{\rm rg}=\chi J_{\rm RP}/\beta_{\rm dS}
\]
therefore remains useful only as a phenomenological probe. It may correlate with \(a_k^2\) or with the stock \(S_k\), but
it is not the fundamental leak law because the healthy receiver itself has no quadratic RP block on the welded family.

\subsection{Approximate split condition from the odd-stock leak law}
\label{subsec:split-threshold}

Let \(R_k^{\rm abs}:=|R_k|\). In the micro-step regime the geometry-side surrogate gives
\[
\Delta E_{{\rm geom},k}
\approx
\frac{\zeta\,(A+c_\Lambda \Lambda_k^2)}{2\sqrt{12\,R_k^{\rm abs}}}\;\Theta_k
\;-\;
\frac{2\zeta\,c_\Lambda \Lambda_k\sqrt{R_k^{\rm abs}}}{\sqrt{12}}\;\delta\Lambda,
\]
while
\[
J_{{\rm wall},k}\approx \mu_{\rm RG}(\Lambda_k)\,\delta\Lambda,
\qquad
\mu_{\rm RG}(\Lambda)\simeq \mu_0+\mu_2\Lambda^2>0.
\]
Hence the local one-wall threshold remains
\begin{equation}
Q_{{\rm leak},k}
\gtrsim
\frac{\zeta\,(A+c_\Lambda \Lambda_k^2)}{2\sqrt{12\,R_k^{\rm abs}}}\;\Theta_k.
\label{eq:split-threshold-base}
\end{equation}
Since
\[
\Theta_k
=
K(\Lambda_k)\,\Delta I_k
\sim
K_0\,\Lambda_k^{-p}\,e^{-2d_{\rm sep}(\Lambda_k)\Lambda_k}\,\Delta I_k,
\qquad p\in\{10,14\},
\]
substituting \eqref{eq:Qleak-effective} gives
\begin{equation}
\widetilde\alpha'(\Lambda_k)\,\delta\Lambda_k\,S_k
\;+\;
\widetilde\alpha(\Lambda_k)\,\eta_k\,g_{\rm cl}(k)\,\Delta I_k
\gtrsim
\frac{\zeta\,(A+c_\Lambda \Lambda_k^2)}{2\sqrt{12\,R_k^{\rm abs}}}\,
K(\Lambda_k)\,\Delta I_k.
\label{eq:split-threshold-effective}
\end{equation}
This is the corrected split threshold. It shows the two late-time branches transparently:
\begin{enumerate}[leftmargin=2em]
\item if \(S_k\) remains negligible, only the fresh term remains, and the exhaustion branch can survive;
\item if \(S_k\) has accumulated, the transport term can dominate even when \(\Delta I_k\) is already tiny.
\end{enumerate}
Under \cref{ass:uv-match} and \(\alpha(\Lambda)\sim \kappa_d\Lambda^{d+2}H^{-d}\), the transport term scales like
\[
\widetilde\alpha'(\Lambda_k)\,\delta\Lambda_k\,S_k
\asymp
\Lambda_k^{d+1}\,\delta\Lambda_k\,S_k,
\]
which is exactly the mechanism by which late split can beat the shrinking fresh payload.

\subsection{Nondegeneracy clause for the normalized seam coefficient}
\label{subsec:separate-c-nondegeneracy}

The split threshold \eqref{eq:split-threshold-effective} is a statement about when the late-time
odd stock can overcome the shrinking fresh payload.
It is \emph{not} by itself a theorem that the normalized seam coefficient \(c_{\Lambda_k}\) on the
canonical compressed welded LoG line is nonzero.
Any theorem proving \(c_{\Lambda_k}\neq 0\) must therefore come from a separate overlap or transport
hypothesis, not from the threshold inequality alone.
Downstream, this clause is absorbed into the delayed selected-line package rather than carried as a
separate package summand.

\begin{definition}[Unit-source seam overlap scalar and no-destructive-cancellation margin]
\label{def:unit-source-overlap-margin}
Fix an accepted tick \(k\).
Assume the canonical compressed welded LoG line is defined and let
\[
\tau^{\mathrm{unit}}_{\Lambda_k}
\]
be a detector-factorized unit welded odd source.
Define the raw unit-source seam overlap scalar by
\begin{equation}
\widetilde c_{\Lambda_k}
:=
\big\langle\!\big\langle
\Pi_{\mathrm{phys},\Lambda_k}\mathsf T_{\Lambda_k}\!\big[\tau^{\mathrm{unit}}_{\Lambda_k}\big],\,
\delta C^{\mathrm{LoG}}_{\Lambda_k}
\big\rangle\!\big\rangle_{C_{\ast,\Lambda_k}}.
\label{eq:unit-source-overlap-scalar}
\end{equation}
A \emph{no-destructive-cancellation margin} at \(\Lambda_k\) is any real number
\[
\mu_{\Lambda_k}>0
\]
such that
\begin{equation}
\abs{\widetilde c_{\Lambda_k}}\ge \mu_{\Lambda_k}.
\label{eq:no-destr-cancel-margin}
\end{equation}
On a connected accepted-branch segment \(I\), a \emph{branchwise no-destructive-cancellation
margin} is a constant \(\mu_I>0\) such that
\begin{equation}
\abs{\widetilde c_{\Lambda}}\ge \mu_I
\qquad
\forall\,\Lambda\in I.
\label{eq:no-destr-cancel-margin-branch}
\end{equation}
\end{definition}

\begin{theorem}[Nondegeneracy criterion for \texorpdfstring{\(c_{\Lambda_k}\)}{c_Lambda_k}]
\label{thm:separate-c-nondegeneracy}
Fix an accepted tick \(k\).
Assume:
\begin{enumerate}[leftmargin=2em]
\item a detector-factorized unit welded odd source at \(\Lambda_k\);
\item the unit-source seam-faithfulness package, namely:
  \begin{enumerate}[label=(\alph*),leftmargin=2em]
  \item the raw seam pairing factorizes as
    \begin{equation}
    \widetilde s_{\Lambda_k}(f,y)
    =
    \sigma_{\mathrm{fill}}(f)\,
    \widetilde c_{\Lambda_k}\,
    \overline{\mathrm{obs}}(y)
    \qquad
    \forall f,\forall y;
    \label{eq:unit-source-seam-factorization}
    \end{equation}
  \item whenever
    \(
    \widetilde c_{\Lambda_k}\neq 0
    \)
    and
    \(
    \overline{\mathrm{obs}}(y)\neq 0
    \),
    the compressed welded LoG line is transversely physical, dominant-line realization holds, and
    the normalized seam coefficient
    \begin{equation}
    c_{\Lambda_k}
    :=
    \kappa_{\Lambda_k}^{-1/2}\,\widetilde c_{\Lambda_k}
    \label{eq:unit-source-normalized-c}
    \end{equation}
    is defined;
  \end{enumerate}
\item a no-destructive-cancellation margin at \(\Lambda_k\) in the sense of
  \cref{def:unit-source-overlap-margin};
\item a witness \(y\) with
  \[
  \overline{\mathrm{obs}}(y)\neq 0.
  \]
\end{enumerate}
Then:
\begin{enumerate}[leftmargin=2em]
\item the raw seam pairing is nonzero for every filler tag \(f\):
  \[
  \widetilde s_{\Lambda_k}(f,y)\neq 0;
  \]
\item the canonical compressed welded LoG line is transversely physical at \(\Lambda_k\);
\item dominant-line realization holds at \(\Lambda_k\);
\item the normalized seam coefficient is defined and nonzero:
  \[
  c_{\Lambda_k}\neq 0.
  \]
\end{enumerate}
If, in addition, for some filler tag \(f\) the source \(\tau^{f,y}_{\Lambda_k}\) is representable on
the characterization layer and admits a source-compatible descendant realization, then the
corresponding descendant representative satisfies
\[
\rho^{\mathrm{dom}}_{\Lambda_k}\!\big(J^{f,y}_{\Lambda_k}\big)\neq 0.
\]
\end{theorem}

\begin{proof}
The margin \eqref{eq:no-destr-cancel-margin} implies
\(
\widetilde c_{\Lambda_k}\neq 0
\).
Together with
\(
\overline{\mathrm{obs}}(y)\neq 0
\),
the factorization \eqref{eq:unit-source-seam-factorization} gives
\(
\widetilde s_{\Lambda_k}(f,y)\neq 0
\)
for every filler tag \(f\), proving item (1).
Item (2), item (3), and the definition and nonvanishing of
\(
c_{\Lambda_k}
\)
are exactly the remaining conclusions of the assumed unit-source seam-faithfulness package under the
same nondegeneracy hypothesis.
The final sentence is the corresponding witness-to-descendant comparison step under representability
and source-compatible descendant realization.
\end{proof}

\begin{corollary}[Branchwise sign-stable nondegeneracy]
\label{cor:branchwise-c-nondegeneracy}
Let \(I\) be a connected accepted-branch segment.
Assume:
\begin{enumerate}[leftmargin=2em]
\item the branchwise normalized seam coefficient
  \[
  \Lambda\mapsto c_{\Lambda}
  \]
  is continuous on \(I\);
\item a branchwise no-destructive-cancellation margin on \(I\) in the sense of
  \cref{def:unit-source-overlap-margin}.
\end{enumerate}
Then
\[
c_{\Lambda}\neq 0
\qquad
\forall\,\Lambda\in I.
\]
In particular, the sign of \(c_{\Lambda}\) is constant on \(I\).
If there also exist a fixed LS witness \(y\) with
\(
\overline{\mathrm{obs}}(y)\neq 0
\)
and a branchwise source-compatible descendant realization satisfying
\[
\rho^{\mathrm{dom}}_{\Lambda}\!\big(J^{f,y}_{\Lambda}\big)
=
\sigma_{\mathrm{fill}}(f)\,
c_{\Lambda}\,
\overline{\mathrm{obs}}(y),
\]
then the sign of
\(
\rho^{\mathrm{dom}}_{\Lambda}(J^{f,y}_{\Lambda})
\)
is constant on \(I\) as well.
\end{corollary}

\begin{proof}
By \eqref{eq:no-destr-cancel-margin-branch}, the scalar
\(
\widetilde c_{\Lambda}
\)
never vanishes on \(I\).
Hence neither does
\(
c_{\Lambda}
\),
because the normalization only rescales the raw overlap scalar by a positive branchwise factor.
A continuous nonzero real-valued function on a connected segment cannot change sign, proving the
constancy of \(\operatorname{sgn}(c_{\Lambda})\).
The final statement follows immediately from the displayed branchwise comparison identity and the
fixed nonzero scalar \(\overline{\mathrm{obs}}(y)\).
\end{proof}

\begin{remark}[Logical separation from the split threshold]
\label{rem:separate-c-nondegeneracy}
The hypotheses of \cref{thm:separate-c-nondegeneracy,cor:branchwise-c-nondegeneracy} are
intentionally separate from \eqref{eq:split-threshold-effective}.
The threshold governs the emergence of split-side stock; the no-destructive-cancellation margin
governs overlap of the unit welded odd source with the canonical compressed dominant line.
Accordingly, the present note treats \(c_{\Lambda}\neq 0\) as the transport/overlap clause later
imported into the delayed selected-line package rather than as a consequence of the split
threshold.
\end{remark}
%
  }{%
    \section{Completion-sector leak model and split criterion}
\label{sec:resolution}

This section separates what is directly derived from UOM from the additional hypotheses required to turn those ingredients
into a concrete split-driving leak model. The new organizing principle is:
the primary hidden variable is a \emph{signed odd split amplitude} \(a_k\), while the positive stock
\(S_k:=a_k^2\) is the quantity that can enter scalar observables and leak laws.
We reserve \(R_k^{\rm abs}\) for the remaining curvature budget and
\[
\mathfrak r_k:=\min_{\Lambda\ge \Lambda_k} f_k(\Lambda)
\]
for the minimized one-wall residual.

\subsection{Odd split mode, output parity, and the rank-one normal form}
\label{subsec:leak-solution}

The only ``leak'' that directly competes with wall work in the \(\Lambda\)-solver is the open-system modification of the
wall energy balance allowed explicitly in UOM's SK/HJ wall equation:
\begin{equation}
\Delta E_{{\rm geom},k}\ =\ J_{{\rm wall},k}\ +\ Q_{{\rm leak},k},
\qquad Q_{{\rm leak},k}\ge 0,
\label{eq:resolution-wall-leak}
\end{equation}
This is the welded wall-balance form used in the earlier summability analysis.
This is a \emph{receiver-side energy complement} term. It must not be conflated with UOM's information-theoretic
bridge/leak instrument from the UOM main manuscript.

\paragraph{What is already derived in UOM.}
Four ingredients are fixed by the manuscript:
\begin{enumerate}[leftmargin=2em]
\item \textbf{Wall balance.} The reduced one-wall description may contain a nonnegative complement \(Q_{\rm leak}\) as in
\eqref{eq:resolution-wall-leak}.
\item \textbf{Quadratic welded SK data.} In the Gaussian sector,
\[
W_\Lambda[J]=\frac12\int J\,(G_{R,\Lambda}+i\,\mathcal N_\Lambda)\,J+O(J^3),
\qquad \mathcal N_\Lambda\ge 0.
\]
\item \textbf{Healthy receiver.} The active KR receiver has vanishing welded RP block, so a receiver-tangent RP penalty is
not the fundamental home of split-driving oddness.
\item \textbf{Continued RP defect.} The analytic-continued covariance has a negative OS mode with
\[
\lambda_{\min}(\mu_{ps},\Lambda)
=
-\alpha(\Lambda)\mu_{ps}^2+O(\mu_{ps}^4),
\qquad
\alpha(\Lambda)\sim \kappa_d\,\Lambda^{d+2}H^{-d}.
\]
\end{enumerate}

\begin{assumption}[Split-mode observability]
\label{ass:wall-match}
Near split onset there exists a one-dimensional odd line \(L_k\) with local coordinate \(a_k\in\mathbb R\) such that the
split involution acts by
\[
a_k\longmapsto -a_k.
\]
The true two-wall advantage
\[
\Delta F_k:=\Delta\mathcal F_{2{\rm w}-1{\rm w},k}
\]
and the minimized one-wall residual
\[
\mathfrak r_k:=\min_{\Lambda\ge \Lambda_k} f_k(\Lambda)
\]
are smooth even scalar outputs of this same local mode.
Equivalently, the 2T \(\Gamma\)-odd descendant direction and the UOM continued negative mode \(u_-(\Lambda_k)\) are
assumed to represent the same active split mode up to a smooth local identification.
\end{assumption}

\begin{lemma}[Evenness of the minimized residual]
\label{lem:residual-even}
Suppose the pre-minimized one-wall residual is a smooth function \(f_k(\Lambda,a)\) satisfying
\[
f_k(\Lambda,-a)=f_k(\Lambda,a),
\]
and suppose there is a unique local minimizer \(\Lambda_\ast(a)\) near \(a=0\) such that
\[
\partial_\Lambda f_k(\Lambda_\ast(a),a)=0,
\qquad
\partial_\Lambda^2 f_k(\Lambda_\ast(0),0)>0.
\]
Then
\[
\mathfrak r_k(a):=f_k(\Lambda_\ast(a),a)
\]
is a smooth even function of \(a\).
\end{lemma}

\begin{proof}
By the implicit function theorem, \(\Lambda_\ast(a)\) is smooth near \(a=0\).
Since \(f_k(\Lambda,a)\) is even in \(a\), the stationarity equation for \(a\) and \(-a\) is identical; uniqueness of the
local minimizer gives
\[
\Lambda_\ast(-a)=\Lambda_\ast(a).
\]
Therefore
\[
\mathfrak r_k(-a)=f_k(\Lambda_\ast(-a),-a)=f_k(\Lambda_\ast(a),a)=\mathfrak r_k(a),
\]
so \(\mathfrak r_k\) is smooth and even.
\end{proof}

\begin{proposition}[Rank-one normal form near split onset]
\label{prop:rank-one-wall-match}
Under \cref{ass:wall-match}, the two scalar outputs admit local even expansions
\begin{equation}
\Delta F_k(a)
=
\mu_k^{(2{\rm w})}+h_k a^2+O(a^4),
\label{eq:DeltaF-detuned}
\end{equation}
\begin{equation}
\mathfrak r_k(a)
=
\mu_k^{(1{\rm w})}+m_k a^2+O(a^4),
\label{eq:rres-detuned}
\end{equation}
with \(h_k,m_k>0\) on the active split line.
\end{proposition}

\begin{proof}
Smooth even functions of one real variable have Taylor series containing only even powers.
The constants \(\mu_k^{(2{\rm w})}\) and \(\mu_k^{(1{\rm w})}\) are the baseline detunings of the two-wall and one-wall
channels. Positivity of \(h_k\) and \(m_k\) expresses that the selected mode is active in both outputs.
\end{proof}

\begin{corollary}[Onset wall-match in rank one]
\label{cor:rank-one-wall-match}
On the onset manifold,
\[
\mu_k^{(1{\rm w})}=\mu_k^{(2{\rm w})}=0,
\]
so
\[
\Delta F_k=h_k a_k^2+O(a_k^4),
\qquad
\mathfrak r_k=m_k a_k^2+O(a_k^4).
\]
Eliminating \(a_k^2\) gives
\begin{equation}
\Delta F_k
=
\frac{h_k}{m_k}\,\mathfrak r_k+O(\mathfrak r_k^2).
\label{eq:rank-one-wall-match}
\end{equation}
\end{corollary}

\subsection{Low-rank generalization}
\label{subsec:alpha-derived}

If the active split sector is \(m\)-dimensional, with odd amplitude vector \(a_k\in\mathbb R^m\), parity gives quadratic
normal forms
\begin{equation}
\Delta F_k(a)
=
\mu_k^{(2{\rm w})}+a^\top H_k a+O(\|a\|^4),
\label{eq:DeltaF-lowrank}
\end{equation}
\begin{equation}
\mathfrak r_k(a)
=
\mu_k^{(1{\rm w})}+a^\top M_k a+O(\|a\|^4),
\label{eq:rres-lowrank}
\end{equation}
with \(H_k,M_k\succeq 0\) on the active split subspace.

\begin{assumption}[Low-rank comparability]
\label{ass:low-rank-comp}
On the active split subspace there exist \(c_k^-,c_k^+>0\) such that
\[
c_k^- M_k \preceq H_k \preceq c_k^+ M_k.
\]
\end{assumption}

\begin{proposition}[Low-rank wall-match comparability]
\label{prop:low-rank-wall-match}
Assume \cref{ass:wall-match,ass:low-rank-comp} and restrict to the onset manifold
\(\mu_k^{(1{\rm w})}=\mu_k^{(2{\rm w})}=0\).
Then for sufficiently small \(\mathfrak r_k\) there exists \(C_k>0\) such that
\begin{equation}
c_k^-\,\mathfrak r_k-C_k\mathfrak r_k^2
\le
\Delta F_k
\le
c_k^+\,\mathfrak r_k+C_k\mathfrak r_k^2.
\label{eq:low-rank-wall-match}
\end{equation}
\end{proposition}

\begin{proof}
On onset,
\[
\Delta F_k=a_k^\top H_k a_k+O(\|a_k\|^4),
\qquad
\mathfrak r_k=a_k^\top M_k a_k+O(\|a_k\|^4).
\]
The quadratic-form comparison gives
\[
c_k^- a_k^\top M_k a_k\le a_k^\top H_k a_k\le c_k^+ a_k^\top M_k a_k.
\]
Since \(\mathfrak r_k\) is itself quadratic in \(a_k\) to leading order, the quartic remainders are \(O(\mathfrak r_k^2)\),
which yields \eqref{eq:low-rank-wall-match}.
\end{proof}

\subsection{Gaussian realization of the odd mode and the leak stock law}
\label{subsec:B-gaussian}

The theorem path above does \emph{not} rely on a KMS-subtracted matrix.
In particular, we do not use
\[
\Sigma_{B,k}:=\bigl(C_{B,k}-C^{\rm KMS}_{\Lambda_k}\bigr)+m_{B,k}m_{B,k}^\top
\]
as a positive object, because it need not be positive semidefinite.
When a safe matrix observable is needed in the realization layer, the natural PSD quantity is
\begin{equation}
M_{B,k}:=C_{B,k}+m_{B,k}m_{B,k}^\top\succeq 0.
\label{eq:MB-def}
\end{equation}

\begin{assumption}[Gaussian completion realization]
\label{ass:gaussian-completion}
For each admissible cutoff \(\Lambda\) in the welded regime, the quadratic influence functional \(W_\Lambda[J]\) admits a
Gaussian dilation, unique up to symplectic equivalence, by a bosonic completion sector \(B_\Lambda\) with canonical
quadratures \(X_\Lambda\) and positive quadratic Hamiltonian
\[
H_{B,\Lambda}=\frac12 X_\Lambda^\top h_\Lambda X_\Lambda,
\qquad h_\Lambda\succeq 0.
\]
\end{assumption}

\begin{assumption}[Dominant odd-mode projection]
\label{ass:export-selection}
Inside the Gaussian completion of \cref{ass:gaussian-completion}, there exists a smooth unit vector
\[
u_k:=u_-(\Lambda_k)
\]
spanning the active odd line and separated from the orthogonal sector by a local spectral gap.
We write
\[
P_k^\perp:=I-u_k u_k^\top
\]
for the orthogonal projector.
\end{assumption}

\begin{assumption}[Classical-control gate]
\label{ass:classical-gate}
There exists a scalar gate \(g_{\rm cl}(k)\in[0,1]\), determined by the accessible classicality defect
\(\mathrm{cl}_k\) or \(\mathrm{Off}_k\), such that \(g_{\rm cl}(k)\ll 1\) at a fresh-aeon start and
\(g_{\rm cl}(k)\to 1\) late in the aeon.
\end{assumption}

\begin{proposition}[Projected Gaussian affine dynamics]
\label{prop:projected-affine}
Under \cref{ass:gaussian-completion,ass:export-selection}, one may choose a finite effective completion-mode basis in which
the completion quadratures obey a linear-Gaussian recursion
\begin{equation}
X_{k+1}=A_k X_k+B_k j_k^{\rm ctl}+\nu_k,
\label{eq:X-recursion}
\end{equation}
with
\[
\mathbb E[\nu_k\mid\mathcal F_k]=0,
\qquad
\mathbb E[\nu_k\nu_k^\top\mid\mathcal F_k]=Q_k\succeq 0.
\]
Consequently the PSD second moment
\[
M_{B,k}=C_{B,k}+m_{B,k}m_{B,k}^\top
\]
obeys
\begin{equation}
M_{B,k+1}=A_k M_{B,k}A_k^\top+Q_k+\text{drift terms}.
\label{eq:MB-recursion}
\end{equation}
Projecting onto the odd line via
\[
a_k:=u_k^\top X_k
\]
gives the scalar affine law
\begin{equation}
a_{k+1}
=
\rho_k a_k+\beta_k+\xi_k+\varepsilon_k^{\rm mix},
\label{eq:B-recursion-general}
\end{equation}
where
\[
\rho_k:=u_{k+1}^\top A_k u_k,\qquad
\beta_k:=u_{k+1}^\top B_k j_k^{\rm ctl},
\]
\[
\xi_k:=u_{k+1}^\top \nu_k,\qquad
\varepsilon_k^{\rm mix}:=u_{k+1}^\top A_k P_k^\perp X_k.
\]
Moreover,
\[
\mathbb E[\xi_k\mid\mathcal F_k]=0,
\qquad
\mathbb E[\xi_k^2\mid\mathcal F_k]=u_{k+1}^\top Q_k u_{k+1}\ge 0.
\]
If \cref{ass:classical-gate} also holds and the welded shape family is fixed, then
\begin{equation}
\mathbb E[\xi_k^2\mid\mathcal F_k]
=
\eta_k\,g_{\rm cl}(k)\,\Delta I_k+O(\Delta I_k^2,\delta\Lambda_k^2),
\label{eq:Zk-linear}
\end{equation}
with \(\eta_k\ge 0\).
\end{proposition}

\begin{proof}
The quadratic SK influence functional determines a Gaussian CPTP map on the finite active completion-mode sector, hence a
linear-Gaussian state-space realization of the form \eqref{eq:X-recursion}; \eqref{eq:MB-recursion} is the standard update
for the PSD second moment.
Decompose \(X_k=u_k a_k+P_k^\perp X_k\), apply \(\ip{u_{k+1}}{\cdot}\) to \eqref{eq:X-recursion}, and collect the projected
transport, deterministic drive, noise, and orthogonal-mode contamination to obtain \eqref{eq:B-recursion-general}.
The mean-zero and positive-variance statements for \(\xi_k\) follow directly from the projected Gaussian noise covariance
\(
u_{k+1}^\top Q_k u_{k+1}\ge 0
\).
Finally, at fixed welded shape family one has \(A_{{\rm eff},k}^2\propto \Delta I_k\), so the fresh covariance injection
into the odd mode is linear in the current payload at leading order; the gate \(g_{\rm cl}(k)\) suppresses this injection
near a fresh-aeon start and yields \eqref{eq:Zk-linear}.
\end{proof}

\begin{assumption}[Centered odd preparation]
\label{ass:centered-odd}
Near split onset, the active odd mode is prepared either branch-symmetrically or in a squeeze-dominated regime, so the
projected deterministic displacement vanishes and orthogonal-mode contamination is higher order:
\[
\beta_k=0,
\qquad
\varepsilon_k^{\rm mix}=O_2.
\]
\end{assumption}

\begin{corollary}[Centered odd-mode law]
\label{cor:centered-odd-mode}
Under
\cref{ass:gaussian-completion,ass:export-selection,ass:classical-gate,ass:centered-odd},
the projected affine law \eqref{eq:B-recursion-general} reduces to
\begin{equation}
a_{k+1}=\rho_k a_k+\xi_k+O_2,
\label{eq:B-recursion}
\end{equation}
with \(\mathbb E[\xi_k\mid\mathcal F_k]=0\) and variance scaling \eqref{eq:Zk-linear}.
\end{corollary}

\paragraph{Stochastic/Gaussian realization layer.}
At theorem level the positive stock is simply
\[
S_k:=a_k^2.
\]
The general projected Gaussian dynamics need not be centered.
Under \cref{ass:centered-odd}, however, the physically relevant late-time regime is the one in which the odd mode carries no
leading deterministic displacement and the fresh effect enters through its variance.
In that centered specialization it is natural to use the conditional second moment
\[
S_k:=\mathbb E[a_k^2\mid\mathcal F_k].
\]

\begin{proposition}[Odd-mode stock law]
\label{prop:odd-mode-stock}
Under
\cref{ass:gaussian-completion,ass:export-selection,ass:classical-gate,ass:centered-odd},
the conditional stock
\[
S_k:=\mathbb E[a_k^2\mid\mathcal F_k]
\]
satisfies
\begin{equation}
S_{k+1}
=
t_k\,S_k+\eta_k\,g_{\rm cl}(k)\,\Delta I_k
+O(S_k\Delta I_k,\Delta I_k^2,\delta\Lambda_k^2),
\qquad t_k=\rho_k^2.
\label{eq:Sk-update}
\end{equation}
\end{proposition}

\begin{proof}
Square \eqref{eq:B-recursion}:
\[
a_{k+1}^2=\rho_k^2 a_k^2+2\rho_k a_k\xi_k+\xi_k^2+O_2.
\]
Taking conditional expectation and using \(\mathbb E[\xi_k\mid\mathcal F_k]=0\) removes the mixed term.
Substituting \eqref{eq:Zk-linear} gives \eqref{eq:Sk-update}.
\end{proof}

\begin{assumption}[UV comparability of energetic stiffness]
\label{ass:uv-match}
There exist constants \(c_-,c_+>0\) and a monotone increasing energetic stiffness \(\widetilde\alpha(\Lambda)\) of the
active odd mode such that, for all sufficiently large \(\Lambda\),
\begin{equation}
c_-\,\alpha(\Lambda)\le \widetilde\alpha(\Lambda)\le c_+\,\alpha(\Lambda).
\label{eq:alpha-comparable}
\end{equation}
\end{assumption}

This is the precise place where the note passes from what is derived in UOM to a physical hypothesis.
The coefficient \(\alpha(\Lambda)\) is derived from the continued OS problem; the statement that the energetic stiffness of
the split mode is comparable to it is not yet a theorem.

With the stock \(S_k\) as the actual state variable, the positive split-driving export flux is modeled by
\begin{equation}
Q_{{\rm leak},k}^{\rm model}
=
\widetilde\alpha'(\Lambda_k)\,\delta\Lambda_k\,S_k
\;+\;
\widetilde\alpha(\Lambda_k)\,\eta_k\,g_{\rm cl}(k)\,\Delta I_k
\;+\; O_2,
\label{eq:Qleak-effective}
\end{equation}
where \(\delta\Lambda_k:=\Lambda_{k+1}-\Lambda_k\).
The first term is the transport cost of carrying already-prepared odd stock to larger cutoff; the second is the fresh
injection due to the current payload.
Under \cref{ass:uv-match}, this gives the scaling comparison
\begin{equation}
Q_{{\rm leak},k}^{\rm model}
\asymp
\alpha'(\Lambda_k)\,\delta\Lambda_k\,S_k
\;+\;
\alpha(\Lambda_k)\,\eta_k\,g_{\rm cl}(k)\,\Delta I_k.
\label{eq:Qleak-alpha}
\end{equation}

\paragraph{Current RP proxy reinterpreted.}
The quantity
\[
Q_{\rm rg}=\chi J_{\rm RP}/\beta_{\rm dS}
\]
therefore remains useful only as a phenomenological probe. It may correlate with \(a_k^2\) or with the stock \(S_k\), but
it is not the fundamental leak law because the healthy receiver itself has no quadratic RP block on the welded family.

\subsection{Approximate split condition from the odd-stock leak law}
\label{subsec:split-threshold}

Let \(R_k^{\rm abs}:=|R_k|\). In the micro-step regime the geometry-side surrogate gives
\[
\Delta E_{{\rm geom},k}
\approx
\frac{\zeta\,(A+c_\Lambda \Lambda_k^2)}{2\sqrt{12\,R_k^{\rm abs}}}\;\Theta_k
\;-\;
\frac{2\zeta\,c_\Lambda \Lambda_k\sqrt{R_k^{\rm abs}}}{\sqrt{12}}\;\delta\Lambda,
\]
while
\[
J_{{\rm wall},k}\approx \mu_{\rm RG}(\Lambda_k)\,\delta\Lambda,
\qquad
\mu_{\rm RG}(\Lambda)\simeq \mu_0+\mu_2\Lambda^2>0.
\]
Hence the local one-wall threshold remains
\begin{equation}
Q_{{\rm leak},k}
\gtrsim
\frac{\zeta\,(A+c_\Lambda \Lambda_k^2)}{2\sqrt{12\,R_k^{\rm abs}}}\;\Theta_k.
\label{eq:split-threshold-base}
\end{equation}
Since
\[
\Theta_k
=
K(\Lambda_k)\,\Delta I_k
\sim
K_0\,\Lambda_k^{-p}\,e^{-2d_{\rm sep}(\Lambda_k)\Lambda_k}\,\Delta I_k,
\qquad p\in\{10,14\},
\]
substituting \eqref{eq:Qleak-effective} gives
\begin{equation}
\widetilde\alpha'(\Lambda_k)\,\delta\Lambda_k\,S_k
\;+\;
\widetilde\alpha(\Lambda_k)\,\eta_k\,g_{\rm cl}(k)\,\Delta I_k
\gtrsim
\frac{\zeta\,(A+c_\Lambda \Lambda_k^2)}{2\sqrt{12\,R_k^{\rm abs}}}\,
K(\Lambda_k)\,\Delta I_k.
\label{eq:split-threshold-effective}
\end{equation}
This is the corrected split threshold. It shows the two late-time branches transparently:
\begin{enumerate}[leftmargin=2em]
\item if \(S_k\) remains negligible, only the fresh term remains, and the exhaustion branch can survive;
\item if \(S_k\) has accumulated, the transport term can dominate even when \(\Delta I_k\) is already tiny.
\end{enumerate}
Under \cref{ass:uv-match} and \(\alpha(\Lambda)\sim \kappa_d\Lambda^{d+2}H^{-d}\), the transport term scales like
\[
\widetilde\alpha'(\Lambda_k)\,\delta\Lambda_k\,S_k
\asymp
\Lambda_k^{d+1}\,\delta\Lambda_k\,S_k,
\]
which is exactly the mechanism by which late split can beat the shrinking fresh payload.

\subsection{Nondegeneracy clause for the normalized seam coefficient}
\label{subsec:separate-c-nondegeneracy}

The split threshold \eqref{eq:split-threshold-effective} is a statement about when the late-time
odd stock can overcome the shrinking fresh payload.
It is \emph{not} by itself a theorem that the normalized seam coefficient \(c_{\Lambda_k}\) on the
canonical compressed welded LoG line is nonzero.
Any theorem proving \(c_{\Lambda_k}\neq 0\) must therefore come from a separate overlap or transport
hypothesis, not from the threshold inequality alone.
Downstream, this clause is absorbed into the delayed selected-line package rather than carried as a
separate package summand.

\begin{definition}[Unit-source seam overlap scalar and no-destructive-cancellation margin]
\label{def:unit-source-overlap-margin}
Fix an accepted tick \(k\).
Assume the canonical compressed welded LoG line is defined and let
\[
\tau^{\mathrm{unit}}_{\Lambda_k}
\]
be a detector-factorized unit welded odd source.
Define the raw unit-source seam overlap scalar by
\begin{equation}
\widetilde c_{\Lambda_k}
:=
\big\langle\!\big\langle
\Pi_{\mathrm{phys},\Lambda_k}\mathsf T_{\Lambda_k}\!\big[\tau^{\mathrm{unit}}_{\Lambda_k}\big],\,
\delta C^{\mathrm{LoG}}_{\Lambda_k}
\big\rangle\!\big\rangle_{C_{\ast,\Lambda_k}}.
\label{eq:unit-source-overlap-scalar}
\end{equation}
A \emph{no-destructive-cancellation margin} at \(\Lambda_k\) is any real number
\[
\mu_{\Lambda_k}>0
\]
such that
\begin{equation}
\abs{\widetilde c_{\Lambda_k}}\ge \mu_{\Lambda_k}.
\label{eq:no-destr-cancel-margin}
\end{equation}
On a connected accepted-branch segment \(I\), a \emph{branchwise no-destructive-cancellation
margin} is a constant \(\mu_I>0\) such that
\begin{equation}
\abs{\widetilde c_{\Lambda}}\ge \mu_I
\qquad
\forall\,\Lambda\in I.
\label{eq:no-destr-cancel-margin-branch}
\end{equation}
\end{definition}

\begin{theorem}[Nondegeneracy criterion for \texorpdfstring{\(c_{\Lambda_k}\)}{c_Lambda_k}]
\label{thm:separate-c-nondegeneracy}
Fix an accepted tick \(k\).
Assume:
\begin{enumerate}[leftmargin=2em]
\item a detector-factorized unit welded odd source at \(\Lambda_k\);
\item the unit-source seam-faithfulness package, namely:
  \begin{enumerate}[label=(\alph*),leftmargin=2em]
  \item the raw seam pairing factorizes as
    \begin{equation}
    \widetilde s_{\Lambda_k}(f,y)
    =
    \sigma_{\mathrm{fill}}(f)\,
    \widetilde c_{\Lambda_k}\,
    \overline{\mathrm{obs}}(y)
    \qquad
    \forall f,\forall y;
    \label{eq:unit-source-seam-factorization}
    \end{equation}
  \item whenever
    \(
    \widetilde c_{\Lambda_k}\neq 0
    \)
    and
    \(
    \overline{\mathrm{obs}}(y)\neq 0
    \),
    the compressed welded LoG line is transversely physical, dominant-line realization holds, and
    the normalized seam coefficient
    \begin{equation}
    c_{\Lambda_k}
    :=
    \kappa_{\Lambda_k}^{-1/2}\,\widetilde c_{\Lambda_k}
    \label{eq:unit-source-normalized-c}
    \end{equation}
    is defined;
  \end{enumerate}
\item a no-destructive-cancellation margin at \(\Lambda_k\) in the sense of
  \cref{def:unit-source-overlap-margin};
\item a witness \(y\) with
  \[
  \overline{\mathrm{obs}}(y)\neq 0.
  \]
\end{enumerate}
Then:
\begin{enumerate}[leftmargin=2em]
\item the raw seam pairing is nonzero for every filler tag \(f\):
  \[
  \widetilde s_{\Lambda_k}(f,y)\neq 0;
  \]
\item the canonical compressed welded LoG line is transversely physical at \(\Lambda_k\);
\item dominant-line realization holds at \(\Lambda_k\);
\item the normalized seam coefficient is defined and nonzero:
  \[
  c_{\Lambda_k}\neq 0.
  \]
\end{enumerate}
If, in addition, for some filler tag \(f\) the source \(\tau^{f,y}_{\Lambda_k}\) is representable on
the characterization layer and admits a source-compatible descendant realization, then the
corresponding descendant representative satisfies
\[
\rho^{\mathrm{dom}}_{\Lambda_k}\!\big(J^{f,y}_{\Lambda_k}\big)\neq 0.
\]
\end{theorem}

\begin{proof}
The margin \eqref{eq:no-destr-cancel-margin} implies
\(
\widetilde c_{\Lambda_k}\neq 0
\).
Together with
\(
\overline{\mathrm{obs}}(y)\neq 0
\),
the factorization \eqref{eq:unit-source-seam-factorization} gives
\(
\widetilde s_{\Lambda_k}(f,y)\neq 0
\)
for every filler tag \(f\), proving item (1).
Item (2), item (3), and the definition and nonvanishing of
\(
c_{\Lambda_k}
\)
are exactly the remaining conclusions of the assumed unit-source seam-faithfulness package under the
same nondegeneracy hypothesis.
The final sentence is the corresponding witness-to-descendant comparison step under representability
and source-compatible descendant realization.
\end{proof}

\begin{corollary}[Branchwise sign-stable nondegeneracy]
\label{cor:branchwise-c-nondegeneracy}
Let \(I\) be a connected accepted-branch segment.
Assume:
\begin{enumerate}[leftmargin=2em]
\item the branchwise normalized seam coefficient
  \[
  \Lambda\mapsto c_{\Lambda}
  \]
  is continuous on \(I\);
\item a branchwise no-destructive-cancellation margin on \(I\) in the sense of
  \cref{def:unit-source-overlap-margin}.
\end{enumerate}
Then
\[
c_{\Lambda}\neq 0
\qquad
\forall\,\Lambda\in I.
\]
In particular, the sign of \(c_{\Lambda}\) is constant on \(I\).
If there also exist a fixed LS witness \(y\) with
\(
\overline{\mathrm{obs}}(y)\neq 0
\)
and a branchwise source-compatible descendant realization satisfying
\[
\rho^{\mathrm{dom}}_{\Lambda}\!\big(J^{f,y}_{\Lambda}\big)
=
\sigma_{\mathrm{fill}}(f)\,
c_{\Lambda}\,
\overline{\mathrm{obs}}(y),
\]
then the sign of
\(
\rho^{\mathrm{dom}}_{\Lambda}(J^{f,y}_{\Lambda})
\)
is constant on \(I\) as well.
\end{corollary}

\begin{proof}
By \eqref{eq:no-destr-cancel-margin-branch}, the scalar
\(
\widetilde c_{\Lambda}
\)
never vanishes on \(I\).
Hence neither does
\(
c_{\Lambda}
\),
because the normalization only rescales the raw overlap scalar by a positive branchwise factor.
A continuous nonzero real-valued function on a connected segment cannot change sign, proving the
constancy of \(\operatorname{sgn}(c_{\Lambda})\).
The final statement follows immediately from the displayed branchwise comparison identity and the
fixed nonzero scalar \(\overline{\mathrm{obs}}(y)\).
\end{proof}

\begin{remark}[Logical separation from the split threshold]
\label{rem:separate-c-nondegeneracy}
The hypotheses of \cref{thm:separate-c-nondegeneracy,cor:branchwise-c-nondegeneracy} are
intentionally separate from \eqref{eq:split-threshold-effective}.
The threshold governs the emergence of split-side stock; the no-destructive-cancellation margin
governs overlap of the unit welded odd source with the canonical compressed dominant line.
Accordingly, the present note treats \(c_{\Lambda}\neq 0\) as the transport/overlap clause later
imported into the delayed selected-line package rather than as a consequence of the split
threshold.
\end{remark}
%
  }%
}

\end{document}
